% Options for packages loaded elsewhere
\PassOptionsToPackage{unicode}{hyperref}
\PassOptionsToPackage{hyphens}{url}
\documentclass[
  ignorenonframetext,
]{beamer}
\newif\ifbibliography
\usepackage{pgfpages}
\setbeamertemplate{caption}[numbered]
\setbeamertemplate{caption label separator}{: }
\setbeamercolor{caption name}{fg=normal text.fg}
\beamertemplatenavigationsymbolsempty
% remove section numbering
\setbeamertemplate{part page}{
  \centering
  \begin{beamercolorbox}[sep=16pt,center]{part title}
    \usebeamerfont{part title}\insertpart\par
  \end{beamercolorbox}
}
\setbeamertemplate{section page}{
  \centering
  \begin{beamercolorbox}[sep=12pt,center]{section title}
    \usebeamerfont{section title}\insertsection\par
  \end{beamercolorbox}
}
\setbeamertemplate{subsection page}{
  \centering
  \begin{beamercolorbox}[sep=8pt,center]{subsection title}
    \usebeamerfont{subsection title}\insertsubsection\par
  \end{beamercolorbox}
}
% Prevent slide breaks in the middle of a paragraph
\widowpenalties 1 10000
\raggedbottom
\AtBeginPart{
  \frame{\partpage}
}
\AtBeginSection{
  \ifbibliography
  \else
    \frame{\sectionpage}
  \fi
}
\AtBeginSubsection{
  \frame{\subsectionpage}
}
\usepackage{iftex}
\ifPDFTeX
  \usepackage[T1]{fontenc}
  \usepackage[utf8]{inputenc}
  \usepackage{textcomp} % provide euro and other symbols
\else % if luatex or xetex
  \usepackage{unicode-math} % this also loads fontspec
  \defaultfontfeatures{Scale=MatchLowercase}
  \defaultfontfeatures[\rmfamily]{Ligatures=TeX,Scale=1}
\fi
\usepackage{lmodern}
\ifPDFTeX\else
  % xetex/luatex font selection
\fi
% Use upquote if available, for straight quotes in verbatim environments
\IfFileExists{upquote.sty}{\usepackage{upquote}}{}
\IfFileExists{microtype.sty}{% use microtype if available
  \usepackage[]{microtype}
  \UseMicrotypeSet[protrusion]{basicmath} % disable protrusion for tt fonts
}{}
\makeatletter
\@ifundefined{KOMAClassName}{% if non-KOMA class
  \IfFileExists{parskip.sty}{%
    \usepackage{parskip}
  }{% else
    \setlength{\parindent}{0pt}
    \setlength{\parskip}{6pt plus 2pt minus 1pt}}
}{% if KOMA class
  \KOMAoptions{parskip=half}}
\makeatother
\setlength{\emergencystretch}{3em} % prevent overfull lines
\providecommand{\tightlist}{%
  \setlength{\itemsep}{0pt}\setlength{\parskip}{0pt}}
% Slide layout defaults for warning-clean scientific decks.
\usepackage{etoolbox}
\IfFileExists{xurl.sty}{\usepackage{xurl}}{}
\IfFileExists{seqsplit.sty}{\usepackage{seqsplit}}{\newcommand{\seqsplit}[1]{#1}}
\protected\def\breakseq#1{\seqsplit{#1}}
\protected\def\breaktt#1{\begingroup\ttfamily\seqsplit{#1}\endgroup}
\setlength{\emergencystretch}{6em}
\tolerance=5000
\hbadness=10000
\hfuzz=1pt
\setlength{\tabcolsep}{2pt}
\AtBeginEnvironment{longtable}{\tiny\renewcommand{\arraystretch}{0.86}\setlength{\tabcolsep}{1pt}}
\AtBeginEnvironment{tabular}{\tiny\renewcommand{\arraystretch}{0.86}\setlength{\tabcolsep}{1pt}}

% Natbib and cross-reference fallbacks — slides don't load natbib
% or cleveref, but combined-PDF manuscript prose may emit these
% commands. The fallback renders citations as a bracketed key list
% and cross-references as detokenized labels so slides don't fail on
% undefined control sequences or raw underscores. \providecommand is
% a no-op if packages load later.
\providecommand{\citep}[1]{[#1]}
\providecommand{\citet}[1]{#1}
\providecommand{\citealp}[1]{#1}
\providecommand{\citeauthor}[1]{#1}
\providecommand{\citeyear}[1]{#1}
\providecommand{\cref}[1]{\texttt{\detokenize{#1}}}
\providecommand{\Cref}[1]{\texttt{\detokenize{#1}}}

\newtheorem{proposition}{Proposition}
\newtheorem{hypothesis}{Hypothesis}
\usepackage{bookmark}
\IfFileExists{xurl.sty}{\usepackage{xurl}}{} % add URL line breaks if available
\urlstyle{same}
% fallback for those not using the hyperref driver hyperxmp:
\makeatletter
\@ifundefined{xmpquote}{\newcommand{\xmpquote}[1]{#1}}{}
\makeatother
\hypersetup{
  hidelinks,
  pdfcreator={LaTeX via pandoc}}

\author{\texorpdfstring{}{}}
\date{}

\begin{document}

\section{Cyber-Phenomenology: Presence, Embodiment, and
Mediation}\label{sec:phenomenology}

\begin{frame}[allowframebreaks]{Cyber-Phenomenology: Presence,
Embodiment, and Mediation}
DigiPPPiP changes the phenomenology of shared drawing. In co-present
PPPiP, the partner's mark arrives with posture, breath, hesitation, and
ambient bodily cues. In remote DigiPPPiP, the mark arrives through a
mediated surface. This does not make the practice unreal; it changes how
embodiment is distributed across body, tool, screen, and archive.
Phenomenology of embodiment gives the theoretical background for
treating tools and media as extensions of bodily agency
\breakseq{[@merleauponty2012phenomenology]}, while presence research
distinguishes social richness, mediated immediacy, and shared-space
illusions as design variables rather than guarantees of togetherness
\breakseq{[@lombard1997presence; @biocca1997cyborg; @lee2004presence]}.

Telepresence and re-embodiment are central. A remote shared canvas
carries traces of the partner's motor act, while video, voice, latency,
cursor trails, stroke halos, and replay create additional layers of
presence and absence \breakseq{[@lombard1997presence; @biocca1997cyborg;
@lee2004presence]}. The concept of digitally mediated or marbled
embodiment helps avoid a simplistic hierarchy in which physical
co-presence is always authentic and remote connection is always degraded
\breakseq{[@atuk2024bodiesonline]}. Video-mediated collaborative drawing also
shows that screen action is not merely a visual output; it is part of
how participants make proposals, invite co-participation, and coordinate
who has practical control of the emerging drawing
\breakseq{[@oittinen2025videodrawing]}. A DigiPPPiP system can support social
presence without photorealistic avatars by using lightweight,
privacy-preserving cues: labeled cursors, delayed traces, visible
turn-taking, and explicit tool-control state.

The same cues can fail phenomenologically. A cursor trail can feel
intimate when it makes hesitation visible, but invasive when it makes
every micro-movement inspectable. A replay can help partners revisit a
moment, but it can also convert a fragile exchange into evidence to be
judged later. DigiPPPiP therefore needs a phenomenology of control as
much as a phenomenology of presence: partners should know when they are
visible, when traces persist, and when a mark can be withdrawn.

Privacy is not just a settings panel in this account.
Contextual-integrity theory treats privacy as appropriate information
flow within a social context, while networked-privacy work treats
privacy as a dynamic boundary-management process rather than a simple
public/private switch \breakseq{[@nissenbaum2011contextualprivacy;
@palen2003privacy]}. Communication Privacy Management theory is
especially relevant for partner drawing because the artifact can become
co-owned: one partner's expressive mark may disclose something about
both partners, and later sharing may require negotiated rules rather
than unilateral export \breakseq{[@petronio2020cpm]}. A DigiPPPiP archive is
therefore phenomenologically legitimate only when persistence, replay,
and deletion preserve the dyad's felt control over the trace.

This matters most for asynchronous practice. The gap between
contributions becomes part of the event structure: the absence after a
mark creates expectation; the next mark answers it. The formal appendix
describes this as curiosity, epistemic action, and co-discovery under an
expected-information-gain arc \breakseq{[@sec:formalisms\_appendix]}. The
phenomenological hypothesis is that waiting can intensify, not merely
interrupt, dyadic meaning when the shared artifact remains stable and
intentional. That claim should be tested with presence measures and
qualitative accounts; it should not be generalized from theory alone.
\end{frame}

\end{document}
